\chapter{Models}\label{ch:models}

\section{Overview of the Modeling Approaches}

This chapter introduces two modeling approaches for spatiotemporal fluid dynamics prediction: a deterministic generative model and a flow-matching generative model. Both approaches share a common hierarchical Swin-based Transformer backbone with additional temporal information flow, but differ in their conditioning strategy and prediction formulation. The sharing of the backbone allows for a fair comparison of the two approaches while isolating the effects of the different generative modeling strategies. The deterministic model directly maps past velocity fields to future states, whereas the flow-matching model learns to conditionally flow a probability distribution.

\subsection{Original Swin Backbone}

The shared backbone is based on the original Swin Transformer \citet{Liu2021SwinTA}, which enables hierarchical feature extraction via shifted window self-attention. The shifted window mechanism allows for efficient local attention computation while maintaining the ability to capture long-range dependencies through hierarchical feature aggregation. The backbone consists of multiple stages, each containing several Swin Transformer blocks, with embedding merging layers between stages to reduce spatial resolution while increasing feature dimensionality and receptive field of applied attention mechanisms. The original Swin Transformer is designed for image classification tasks, but we adapt it for spatiotemporal fluid dynamics prediction by incorporating temporal information and modifying the attention mechanism to handle sequential data. Two challenges here are the extra temporal dimension and the need to return the processed data to the original spatial resolution instead of a classification output.

\subsection{Space-Then-Time Attention Backbone}

To incorporate temporal structure, we employ a custom space-then-time attention block. As temporal interactions, spacings and lengths are typically different than spatial interactions, it makes sense to factorize the attention mechanism into two seperate stages. Additionally, using a seperate attention stage for the temporal dimension allows for possibly more efficient and effective modeling of temporal dependencies, which is highly relevnant for large-scale foundation training tasks. This is motivated by the paper >>>>>>>>>>>>>>BCAT, which concludes that a space-then-time mechanism is more efficient while negligibly affecting predictive performance.
In this formulation, self-attention is first applied across spatial dimensions using the spatial windows of the original swin transformer block and subsequently in straight line across the temporal dimension. The corresponding attention masking strategy is illustrated in Figure~\ref{fig:BlockExplainer}. In practice, multiple Space-Then-Time attention blocks are used sequentially for smoother and more robust information spread.

#################### include evidence as temporal bcat model thingy.

\begin{figure}[h]
	\centering
	\makebox[\textwidth][c]{
		\includestandalone[width=0.8\textwidth]{figs/BlockExplainerAll}
	}
	\caption{Illustration of the space-then-time attention masking pattern used in the backbone architecture. The left and middle figure show the two subsequent steps of the spatial attention mask and the right figure shows the temporal attention mask.}
	\label{fig:BlockExplainer}
\end{figure}

\subsection{UNet-style Multi-scale Architecture}

The Swin Transformer backbone is embedded within a UNet-style encoder-decoder structure with skip connections. This design enables multi-scale feature extraction while preserving fine-grained spatial information. Additionally, the output dimensions of the model are the same size as the input, which is not the case with the classification output of the traditional Swin Transformer. Skip connections are implemented using ConvNeXt-style blocks, which facilitate stable gradient propagation and information flow across resolution scales. Each level of the Unet focusses on a different spatial scale of information spread, with the deepest layer programmed for an spatial all-to-all attention mask.

\subsection{Patch Embedding and Reconstruction}

At the input stage, the two-dimensional velocity field is partitioned into non-overlapping square patches. Each patch is projected into a latent embedding space using a learned 2-layer MLP. Static positional and temporal embeddings are added to each patch embedding to encode spatial and temporal structure. A single patch resembles both the x- and y-component of a part of the velocity field for a single time step.

At the output stage, the latent embeddings are projected back to the original velocity field dimension using a learned 3-layer MLP and reassembled into the full spatial grid.

\subsection{Hierarchical Patch Merging and Expansion}

As depth increases in the encoder, groups of four neighbouring patches are merged into a single token with double size embedding dimension with a single dense connected layer. This operation progressively reduces spatial resolution while increasing receptive field size. This form of compression throughout the depth of the model also realises a form of regularization and generalisation. In the decoder, the reverse operation is applied to recover full spatial resolution.

\section{Deterministic Model}

\subsection{Prediction}

The deterministic model receives as input a sequence of the past five velocity field states and directly predicts the next five states. Due to the grouping of five subsequent states of the system, we aim to reduce the amount of compounding error from recursive inference.


\subsection{Architecture Overview}

The deterministic model uses the shared backbone described in the previous section. It operates on the same patchified spatiotemporal representation and produces a direct mapping from past states to subsequent future states.

\begin{figure}[h]
	\centering
	\makebox[\textwidth][c]{
		\includestandalone[width=1\textwidth]{figs/FM}
	}
	\caption{Architecture overview of the deterministic model.}
	\label{fig:det_model_architecture}
\end{figure}

\subsection{Hyperparameter Configurations}

Two deterministic model variants are evaluated. One uses smaller patch sizes with higher embedding dimensionality and greater depth, while the other uses larger patch sizes with reduced embedding dimensionality and depth. This design allows exploration of the trade-off between local resolution and global context modeling.


\begin{table}[th]
\centering
\caption{Hyperparameters for deterministic FluidGPT models.}
\label{tab:fluidgpt_hparams_det}
\begin{tabular}{lcc}
\toprule
\textbf{Hyperparameter} & \textbf{AR-d2} & \textbf{AR-d3} \\
\midrule
Temporal Bundling             & 5            & 5            \\
Embedding Dimension           & 96           & 48           \\
Embedder Encoding/Decoding    & linear       & linear       \\
Patch Shape                   & $8{\times}8{\times}2$ & $4{\times}4{\times}2$ \\
Embedder Hidden Dim.          & 512          & 256          \\
Model Depth                   & 2            & 3            \\
Stage Blocks (per-depth)      & [3,3,3,3,3]  & [2,2,2,2,2,2,2] \\
Number of Heads (per-depth)   & [4,8,16,8,4] & [4,8,8,16,8,8,4] \\
Window Attention Size         & $4{\times}4$ & $4{\times}4$ \\
Activation Function           & GELU         & GELU         \\
Skip Block                    & convnext     & convnext     \\
Gradient Flowthrough Skips    & true (all)   & true (all)   \\
Model Parameters              & 24.7\,M      & 17.9\,M      \\
\bottomrule
\end{tabular}
\end{table}


\section{Flow-Matching Model}

\subsection{Prediction}

The flow-matching model conditions on the same five past velocity field states as the deterministic model. In addition, it concatenates a noise sample drawn from a prior distribution. The model learns to predict a velocity field that transforms the noisy sample into a physically consistent future state. Note that the term velocity does not refer to the physical velocity of the Navier-Stokes flow, but instead refers to the flow from noisy distribution to the learned distribution of subsequent states. Even though the model itself is still deterministic, the nature of sampling from the prior distribution makes this model a probabilistic model.

\subsection{Flow-Time Conditioning}

Flow-time is encoded using a high-dimensional sinusoidal embedding, which is processed by an MLP with SiLU activations. This embedding is reintroduced with concatenation onto the embeddings at each start of the Space-Then-Time attention block sequence. This ensuring that temporal information is available throughout the network and is not lost during the attention processes.

\subsection{Prior Distribution Construction}

The prior distribution is generated by applying a Gaussian filter to pure Gaussian noise. This introduces spatial and temporal correlations, resulting in a structured noise field. The standard deviations for the Gaussian filter are set to 4 in spatial dimensions and 1 in the temporal dimension, based on empirical tuning. Putting a Gaussian filter over the raw Gaussian noise acts as a low-pass filter
resulting in a easier learning problem and thus decreasing the amount of time required for the model to converge while training.

\subsection{Sampling and Integration}

The model predicts a velocity increment field conditioned on flow-time. The final predicted state is obtained by integrating this velocity field over flow-time using a numerical integration scheme. We use a euler-integration scheme with 20 steps.

\subsection{Architecture Overview}

The flow-matching model shares the same backbone architecture as the deterministic model, with the primary difference being the inclusion of flow-time conditioning and noise input.

\begin{figure}[h]
	\centering
	\makebox[\textwidth][c]{
		\includestandalone[width=1\textwidth]{figs/FM}
	}
	\caption{Architecture overview of the flow-matching model.}
	\label{fig:fm_model_architecture}
\end{figure}

\subsection{Hyperparameter Configurations}

As with the deterministic model, two variants are evaluated for flow matching. These variants differ in patch size, embedding dimension, depth, and number of attention blocks per layer, enabling a systematic study of model capacity versus resolution trade-offs.


\begin{table}[H]
\centering
\caption{Hyperparameters for probabilistic FluidGPT models.}
\label{tab:fluidgpt_hparams_prob}
\begin{tabular}{lcc}
\toprule
\textbf{Hyperparameter} & \textbf{FM-d2} & \textbf{FM-d3} \\
\midrule
Temporal Bundling             & 5 + 5 (cond.) & 5 + 5 (cond.) \\
Embedding Dimension           & 128          & 64           \\
Embedder Encoding/Decoding    & linear       & linear       \\
Patch Shape                   & $8{\times}8{\times}2$ & $4{\times}4{\times}2$ \\
Embedder Hidden Dim.          & 512          & 256          \\
Flow-time Embedding Dim.      & 256          & 256          \\
Flow-time Injection Size      & [32,64,128]  & [16,32,64,128] \\
Model Depth                   & 2            & 3            \\
Stage Blocks (per-depth)      & [3,3,3,3,3]  & [2,2,2,2,2,2,2] \\
Number of Heads (per-depth)   & [8,16,32,16,8] & [4,8,16,32,16,8,4] \\
Window Attention Size         & $4{\times}4$ & $4{\times}4$ \\
Attention Masking             & not on cond. & not on cond. \\
Activation Function           & GELU         & GELU         \\
Skip Block                    & convnext     & convnext     \\
Gradient Flowthrough Skips    & true (all)   & true (all)   \\
Model Parameters              & 64.3\,M      & 46.0\,M      \\
\bottomrule
\end{tabular}
\end{table}

\section{Model Scaling and Architectural Variants}

Both modeling approaches are evaluated under two architectural regimes. One configuration uses smaller patch sizes with higher embedding dimensions and deeper networks, while the other uses larger patch sizes with reduced embedding capacity. The size of mlp layers is adjusted accordingly to maintain a consistent computational budget for its patch size. Additionally, the number of space-then-time attention blocks per layer is varied as a hyperparameter to control computational cost and define less difference in representational capacity between the d2 and d3-models. These choices allow us to investigate the trade-off between fine-grained spatial resolution and global contextual reasoning. 

Note that patch fusion results in a double size embedding dimension and fourfold pixel size receptiveness, thus making the second depth embeddings of the d-3 models and the first depth embeddings of the d-2 models equivalent in size and pixel-correspondence.

